% !TeX program = xelatex
\documentclass[12pt,a4paper,oneside]{report}

\usepackage[a4paper,top=2.3cm,bottom=2.4cm,left=2.2cm,right=2.2cm,headheight=40pt]{geometry}
\usepackage{xcolor}
\usepackage{graphicx}
\usepackage{booktabs}
\usepackage{longtable}
\usepackage{tabularx}
\usepackage{array}
\usepackage{multirow}
\usepackage{float}
\usepackage{caption}
\usepackage{subcaption}
\usepackage{enumitem}
\usepackage{ragged2e}
\usepackage{amsmath,amssymb}
\usepackage{tikz}
\usetikzlibrary{arrows.meta,positioning,shapes.geometric,fit,backgrounds,calc}
\usepackage{pgfplots}
\pgfplotsset{compat=1.18}
\usepackage[most]{tcolorbox}
\usepackage{fancyhdr}
\usepackage{titlesec}
\usepackage{listings}
\usepackage{hyperref}
\usepackage{xepersian}

% ---------- Editable metadata ----------
\newcommand{\UniversityName}{شریف صنعتی دانشگاه}
\newcommand{\FacultyName}{کامپیوتر مهندسی دانشکده}
\newcommand{\CourseName}{پروژه هوش مصنوعی}
\newcommand{\InstructorName}{دکتر تن قطاری}
\newcommand{\StudentOne}{فاطمه بهزادی‌ دیل}
\newcommand{\StudentOneID}{403105809}
\newcommand{\StudentTwo}{فاطمه دمیرچی}
\newcommand{\StudentTwoID}{403110241}
\newcommand{\RepoURL}{https://github.com/FDamirchi/sign-language-translator}
\newcommand{\UniversityLogo}{assets/sharif_logo_blue.png}

% ---------- Fonts ----------
% XeLaTeX only. The file first tries locally bundled fonts, then installed fonts,
% and finally open fallback fonts so compilation never stops at font loading.
% Put the two font files in a folder named "fonts" next to this .tex file.
% Common filenames are detected automatically.

\newcommand{\SetupNazaninFromFile}[2]{%
  \settextfont[
    Path={#1},
    UprightFont={#2},
    BoldFont={#2},
    ItalicFont={#2},
    BoldItalicFont={#2},
    Script=Arabic,
    Scale=1.16
  ]{#2}%
  \newfontfamily\MetricFont[
    Path={#1},
    UprightFont={#2},
    BoldFont={#2},
    Script=Arabic,
    Scale=1.16
  ]{#2}%
}

\newcommand{\SetupTitrFromFile}[2]{%
  \newfontfamily\TitleFont[
    Path={#1},
    UprightFont={#2},
    BoldFont={#2},
    Script=Arabic,
    Scale=1.02
  ]{#2}%
}

% Text font: local file -> installed family -> fallback.
\IfFileExists{fonts/BNAZANIN.TTF}{%
  \SetupNazaninFromFile{fonts/}{BNAZANIN.TTF}%
}{\IfFileExists{fonts/BNazanin.ttf}{%
  \SetupNazaninFromFile{fonts/}{BNazanin.ttf}%
}{\IfFileExists{BNAZANIN.TTF}{%
  \SetupNazaninFromFile{}{BNAZANIN.TTF}%
}{\IfFontExistsTF{B Nazanin}{%
  \settextfont[
    BoldFont={B Nazanin},ItalicFont={B Nazanin},BoldItalicFont={B Nazanin},
    Script=Arabic,Scale=1.16
  ]{B Nazanin}%
  \newfontfamily\MetricFont[
    BoldFont={B Nazanin},Script=Arabic,Scale=1.16
  ]{B Nazanin}%
}{%
  \settextfont[Script=Arabic,Scale=1.05]{DejaVu Sans}%
  \newfontfamily\MetricFont[Script=Arabic,Scale=1.05]{DejaVu Sans}%
}}}}

% Heading font: local file -> installed family -> fallback.
\IfFileExists{fonts/BTITRBD.TTF}{%
  \SetupTitrFromFile{fonts/}{BTITRBD.TTF}%
}{\IfFileExists{fonts/BTITR.TTF}{%
  \SetupTitrFromFile{fonts/}{BTITR.TTF}%
}{\IfFileExists{BTITRBD.TTF}{%
  \SetupTitrFromFile{}{BTITRBD.TTF}%
}{\IfFontExistsTF{B Titr}{%
  \newfontfamily\TitleFont[BoldFont={B Titr},Script=Arabic,Scale=1.02]{B Titr}%
}{%
  \newfontfamily\TitleFont[Script=Arabic,Scale=0.95]{DejaVu Sans}%
}}}}

\setlatintextfont[Scale=0.92]{DejaVu Sans}
\setmonofont[Scale=0.84]{DejaVu Sans Mono}

% ---------- Colors inspired by the university identity and AI interfaces ----------
\definecolor{Primary}{HTML}{0B4F8A}
\definecolor{PrimaryDark}{HTML}{082F55}
\definecolor{Secondary}{HTML}{00A6A6}
\definecolor{Accent}{HTML}{6C63FF}
\definecolor{Warm}{HTML}{F59E0B}
\definecolor{SoftBlue}{HTML}{EEF6FC}
\definecolor{SoftGreen}{HTML}{EAF9F7}
\definecolor{SoftPink}{HTML}{F4F0FF}
\definecolor{SoftGray}{HTML}{F4F7FA}
\definecolor{DarkGray}{HTML}{243447}
\definecolor{CodeBg}{HTML}{F7F9FC}

% ---------- General settings ----------
\hypersetup{
  colorlinks=true,
  linkcolor=Primary,
  urlcolor=Accent,
  citecolor=Secondary,
  pdftitle={گزارش پروژه مترجم بلادرنگ زبان اشاره},
  pdfauthor={\StudentOne، \StudentTwo}
}
\setlength{\parindent}{0pt}
\setlength{\parskip}{6pt}
\linespread{1.38}
\AtBeginDocument{\setRTL\justifying}
\setlist[itemize]{label=\textcolor{Secondary}{\scriptsize$\blacktriangleright$},itemsep=3pt,topsep=4pt,rightmargin=1.2em}
\setlist[enumerate]{itemsep=3pt,topsep=4pt,rightmargin=1.2em}
\captionsetup{font=small,labelfont=bf,justification=centering}

\titleformat{\chapter}[display]
  {\TitleFont\fontsize{25}{34}\selectfont\color{Primary}\raggedleft}
  {\TitleFont\color{Secondary}\Large\raggedleft فصل \thechapter}
  {5pt}{{\color{Primary!35}\titlerule[1.1pt]}\vspace{7pt}\color{Primary}\raggedleft}
\titleformat{\section}
  {\TitleFont\fontsize{18}{26}\selectfont\color{Primary}\raggedleft}
  {\thesection}{0.65em}{}
\titleformat{\subsection}
  {\TitleFont\fontsize{15}{22}\selectfont\color{Secondary}\raggedleft}
  {\thesubsection}{0.65em}{}



\pagestyle{fancy}
\fancyhf{}
\fancyhead[R]{\smash{\small\color{DarkGray}گزارش پروژه مترجم بلادرنگ زبان اشاره}}
\fancyhead[L]{\smash{\small\color{DarkGray}\CourseName}}
\fancyfoot[C]{\color{Primary}\thepage}
\renewcommand{\headrulewidth}{0.45pt}
\renewcommand{\headrule}{\hbox to\headwidth{\color{Primary!30}\leaders\hrule height \headrulewidth\hfill}}
\renewcommand{\footrulewidth}{0pt}

\newcolumntype{Y}{>{\RTL\raggedleft\arraybackslash}X}
\newcolumntype{C}[1]{>{\centering\arraybackslash}p{#1}}
\newcolumntype{R}[1]{>{\RTL\raggedleft\arraybackslash}p{#1}}

\newtcolorbox{infobox}[1][]{
  enhanced,breakable,colback=SoftBlue,colframe=Primary,
  boxrule=0.7pt,arc=3mm,left=4mm,right=4mm,top=2.5mm,bottom=2.5mm,
  title=#1,fonttitle=\TitleFont,coltitle=Primary
}
\newtcolorbox{successbox}[1][]{
  enhanced,breakable,colback=SoftGreen,colframe=Secondary,
  boxrule=0.7pt,arc=3mm,left=4mm,right=4mm,top=2.5mm,bottom=2.5mm,
  title=#1,fonttitle=\TitleFont,coltitle=Secondary
}
\newtcolorbox{warningbox}[1][]{
  enhanced,breakable,colback=SoftPink,colframe=Accent,
  boxrule=0.7pt,arc=3mm,left=4mm,right=4mm,top=2.5mm,bottom=2.5mm,
  title=#1,fonttitle=\TitleFont,coltitle=Accent
}
\newtcolorbox{codebox}[1][]{
  enhanced,breakable,colback=CodeBg,colframe=DarkGray,
  boxrule=0.5pt,arc=2mm,left=3mm,right=3mm,top=2mm,bottom=2mm,
  title=#1,fonttitle=\TitleFont,coltitle=DarkGray
}

\lstdefinestyle{pythonstyle}{
  language=Python,
  basicstyle=\ttfamily\footnotesize,
  keywordstyle=\color{Accent}\bfseries,
  commentstyle=\color{Secondary},
  stringstyle=\color{Warm},
  backgroundcolor=\color{CodeBg},
  frame=single,
  rulecolor=\color{DarkGray!35},
  breaklines=true,
  showstringspaces=false,
  numbers=left,
  numberstyle=\tiny\color{DarkGray!70},
  xleftmargin=1.2em,
  framexleftmargin=0.8em,
  tabsize=4
}

% ---------- Helper commands ----------
\newcommand{\eng}[1]{\lr{#1}}
\newcommand{\code}[1]{\lr{\nolinkurl{#1}}}
\newcommand{\metric}[1]{\lr{{\MetricFont\textbf{#1}}}}
\newcommand{\yes}{\textcolor{Secondary}{\bfseries انجام‌شده}}
\newcommand{\partialdone}{\textcolor{Warm}{\bfseries نیازمند تکمیل}}

% Compile even when optional report figures have not been copied next to the .tex file.
% As soon as the referenced file exists, the real image is used automatically.
\newcommand{\safeincludegraphics}[2][]{%
  \IfFileExists{#2}{\includegraphics[#1]{#2}}{%
    \begingroup
    \setlength{\fboxsep}{6pt}%
    \fbox{\parbox[c][3.2cm][c]{0.78\linewidth}{%
      \centering\small\textcolor{DarkGray}{تصویر یافت نشد}\\[3pt]
      \scriptsize\lr{\detokenize{#2}}%
    }}%
    \endgroup
  }%
}

\begin{document}
\setRTL
% Reapply heading styles after the bidi/xepersian patches are loaded
\titleformat{\chapter}[display]
  {\TitleFont\fontsize{25}{35}\selectfont\color{Primary}\raggedleft}
  {\TitleFont\fontsize{17}{24}\selectfont\color{Secondary}\raggedleft فصل \thechapter}
  {5pt}{{\color{Primary!30}\titlerule[1.1pt]}\vspace{9pt}\color{Primary}\raggedleft}

% ==================== Title page ====================
\begin{titlepage}
\begin{tikzpicture}[remember picture,overlay]
  \shade[top color=PrimaryDark,bottom color=Primary] (current page.north west)
    rectangle ([yshift=-6.1cm]current page.north east);
  \fill[Secondary] ([yshift=-6.1cm]current page.north west)
    rectangle ([yshift=-6.22cm]current page.north east);
  \fill[SoftGray] (current page.south west)
    rectangle ([yshift=2.05cm]current page.south east);

  % Minimal neural-network/circuit pattern in the cover background
  \foreach \x/\y in {2.0/-1.1,4.1/-2.2,6.0/-1.25,8.2/-2.7,10.5/-1.4,12.8/-2.35,15.2/-1.15,17.3/-2.55} {
    \fill[white,opacity=0.16] ([xshift=\x cm,yshift=\y cm]current page.north west) circle (1.7pt);
  }
  \draw[white,opacity=0.12,line width=0.7pt]
    ([xshift=2.0cm,yshift=-1.1cm]current page.north west) --
    ([xshift=4.1cm,yshift=-2.2cm]current page.north west) --
    ([xshift=6.0cm,yshift=-1.25cm]current page.north west) --
    ([xshift=8.2cm,yshift=-2.7cm]current page.north west) --
    ([xshift=10.5cm,yshift=-1.4cm]current page.north west) --
    ([xshift=12.8cm,yshift=-2.35cm]current page.north west) --
    ([xshift=15.2cm,yshift=-1.15cm]current page.north west) --
    ([xshift=17.3cm,yshift=-2.55cm]current page.north west);

  \node[fill=white,draw=white,minimum size=2.35cm,circle,inner sep=4mm]
    at ([xshift=-2.35cm,yshift=-2.35cm]current page.north east)
    {\IfFileExists{\UniversityLogo}{\includegraphics[width=1.65cm]{\UniversityLogo}}{\TitleFont\fontsize{20}{24}\selectfont\textcolor{Primary}{ش}}};
  \node[anchor=east,text=white,font=\TitleFont\fontsize{18}{25}\selectfont]
    at ([xshift=-4.0cm,yshift=-1.95cm]current page.north east) {\UniversityName};
  \node[anchor=east,text=white!85,font=\fontsize{12}{18}\selectfont]
    at ([xshift=-4.0cm,yshift=-2.75cm]current page.north east) {\FacultyName};
  \node[anchor=east,text=white!48,font=\fontsize{46}{50}\selectfont]
    at ([xshift=-1.2cm,yshift=-5.05cm]current page.north east) {\lr{AI}};
\end{tikzpicture}

\vspace*{6.55cm}
\begin{flushright}
  {\TitleFont\fontsize{28}{39}\selectfont\textcolor{Primary}{گزارش جامع پروژه شماره ۳}}\par
  \vspace{0.55cm}
  {\TitleFont\fontsize{20}{31}\selectfont\textcolor{DarkGray}{طراحی و پیاده‌سازی سیستم طبقه‌بندی تصاویر زبان اشاره و مترجم بلادرنگ به گفتار}}\par
  \vspace{0.35cm}
  {\color{Secondary}\rule{7.2cm}{1.5pt}}\par
\end{flushright}

\vspace{1.0cm}
\begin{tcolorbox}[
  enhanced,colback=white,colframe=Primary!28,boxrule=0.75pt,arc=4mm,
  shadow={1.5mm}{-1.5mm}{0mm}{black!12},left=7mm,right=7mm,top=5mm,bottom=5mm,
  borderline east={2.4pt}{0pt}{Secondary}
]
\begin{tabularx}{\textwidth}{R{4.3cm}Y}
\textbf{نام درس} & \CourseName \\
\textbf{استاد} & \InstructorName \\
\textbf{دانشجوی اول} & \StudentOne\quad \textcolor{Secondary}{\lr{|}} \quad \StudentOneID \\
\textbf{دانشجوی دوم} & \StudentTwo\quad \textcolor{Secondary}{\lr{|}} \quad \StudentTwoID \\
\textbf{مخزن پروژه} & \href{\RepoURL}{\lr{github.com/FDamirchi/sign-language-translator}}
\end{tabularx}
\end{tcolorbox}

\vfill
\begin{center}
\end{center}
\end{titlepage}

\pagenumbering{alph}
\chapter*{چکیده}
\addcontentsline{toc}{chapter}{چکیده}
این پروژه یک سامانه انتهابه‌انتهای تشخیص حروف زبان اشاره آمریکایی از تصویر، تجمیع خروجی به متن و تبدیل متن نهایی به گفتار را پیاده‌سازی می‌کند. داده مورد استفاده شامل \metric{۸۷\,۰۰۰} تصویر آموزشی در \metric{۲۹} کلاس است: حروف \eng{A--Z} و سه فرمان \eng{del}، \eng{nothing} و \eng{space}. تحلیل اکتشافی نشان می‌دهد تمام کلاس‌های مجموعه آموزشی دقیقاً \metric{۳\,۰۰۰} تصویر دارند، تصاویر همگی \eng{RGB} و با اندازه ثابت \eng{200×200} پیکسل هستند و در نتیجه داده از نظر تعداد نمونه کاملاً متوازن است.

در فاز مدل‌سازی، یک شبکه عصبی پیچشی چهاربخشی با \metric{۱۹\,۲۷۸\,۱۷۳} پارامتر طراحی و با چارچوب \eng{PyTorch} آموزش داده شد. داده‌ها با نسبت‌های \metric{۷۵/۱۵/۱۰} به آموزش، اعتبارسنجی و آزمون تقسیم شدند. مدل با بهینه‌ساز \eng{Adam}، نرخ یادگیری \eng{0.001}، اندازه دسته \eng{64} و به‌مدت \eng{10} دوره آموزش دید. بهترین دقت اعتبارسنجی در دوره نهم برابر \metric{99.23\%} و دقت آزمون نهایی برابر \metric{99.34\%} گزارش شد.

در فاز بلادرنگ، یک معماری ماژولار برای دریافت تصویر وب‌کم، استخراج ناحیه دست، همسان‌سازی پیش‌پردازش با مدل، استنتاج روی \eng{CPU/GPU}، تثبیت زمانی پیش‌بینی، مدیریت فرمان‌های ویرایشی، تشخیص پایان رشته و پخش صوت با \eng{pyttsx3} ایجاد شد. علامت تنها زمانی ثبت می‌شود که با اطمینان حداقل \metric{80\%} برای \metric{۲ ثانیه} پایدار بماند؛ کلاس \eng{NOTHING} پس از \metric{۰.۴ ثانیه} قفل علامت قبلی را آزاد و پس از \metric{۵ ثانیه} متن موجود را نهایی می‌کند. کد با تست‌های واحد برای اجزای اصلی و کنترل کیفیت با \eng{pytest} و \eng{Ruff} همراه شده است.

\tableofcontents
\listoffigures
\listoftables
\clearpage
\pagenumbering{arabic}

% ==================== Chapter 1 ====================
\chapter{تعریف مسئله، اهداف و معیارهای تحویل}

\section{بیان مسئله}
هدف پروژه توسعه یک سیستم کامل هوش مصنوعی از مرحله شناخت داده خام تا استقرار در یک محیط کاربردی بلادرنگ است. ورودی سیستم فریم‌های دوربین و خروجی آن متن حروف‌چینی‌شده و صوت متن نهایی است. مسئله اصلی را می‌توان به سه زیرمسئله تقسیم کرد:
\begin{enumerate}
  \item تحلیل و استانداردسازی تصاویر کلاس‌بندی‌شده زبان اشاره؛
  \item آموزش یک طبقه‌بند \eng{CNN} برای تشخیص یکی از ۲۹ کلاس؛
  \item طراحی منطق زمانی برای تبدیل پیش‌بینی‌های فریم‌به‌فریم به رشته پایدار و قابل خواندن.
\end{enumerate}

\section{اهداف فنی}
\begin{itemize}
  \item بررسی توازن کلاس‌ها، ویژگی‌های تصویر و توزیع داده؛
  \item طراحی شبکه شامل لایه‌های \eng{Convolution}، \eng{MaxPooling} و \eng{Fully Connected}؛
  \item کاهش بیش‌برازش با افزایش داده، \eng{Dropout} و تنظیم نرخ یادگیری؛
  \item دریافت تصویر با \eng{OpenCV} و محدودسازی ورودی مدل به یک ناحیه مربعی؛
  \item ثبت علامت تنها پس از تشخیص پایدار برای جلوگیری از نویز لحظه‌ای؛
  \item پشتیبانی از فرمان‌های \eng{SPACE}، \eng{DEL} و \eng{NOTHING}؛
  \item تبدیل رشته نهایی به گفتار و آماده‌سازی سیستم برای عبارت بعدی؛
  \item جداسازی اجزا و فراهم‌کردن تست‌های واحد برای توسعه و نگه‌داری مطمئن.
\end{itemize}

\section{تطبیق الزامات صورت پروژه با خروجی نهایی}
\begin{table}[H]
\centering
\caption{پوشش الزامات اصلی پروژه}
\begin{tabularx}{\textwidth}{R{3.2cm}Y C{2.5cm}}
\toprule
\textbf{فاز} & \textbf{الزام} & \textbf{وضعیت} \\
\midrule
فاز ۱ & شمارش نمونه‌ها و بررسی توازن ۲۹ کلاس & \yes \\
فاز ۱ & نمایش نمونه تصادفی از هر کلاس در شبکه تصویری & \yes \\
فاز ۱ & تحلیل ابعاد، کانال رنگ و بازه پیکسل‌ها & \yes \\
فاز ۱ & تقسیم داده به آموزش، اعتبارسنجی و آزمون & \yes \\
فاز ۲ & طراحی و آموزش \eng{CNN} و شرح معماری & \yes \\
فاز ۲ & بررسی اثر پارامترهایی نظیر نرخ یادگیری، اندازه دسته و بهینه‌ساز & \yes \\
فاز ۲ & ارزیابی روی داده آزمون و گزارش دقت & \yes \\
فاز ۲ & نمودار \eng{Accuracy} و \eng{Loss} و ماتریس درهم‌ریختگی & \yes \\
فاز ۳ & دریافت تصویر وب‌کم، کادر دست و ارسال \eng{ROI} به مدل & \yes \\
فاز ۳ & تثبیت علامت برای ۲ ثانیه با اطمینان بالا & \yes \\
فاز ۳ & نهایی‌سازی بعد از ۵ ثانیه \eng{NOTHING} & \yes \\
فاز ۳ & تبدیل متن به گفتار و پاک‌کردن رشته & \yes \\
\bottomrule
\end{tabularx}
\end{table}

% ==================== Chapter 2 ====================
\chapter{فاز اول: تحلیل اکتشافی و پیش‌پردازش داده}

\section{ساختار مجموعه‌داده}
مجموعه آموزشی شامل ۲۹ پوشه کلاس است. ۲۶ کلاس نخست حروف الفبای انگلیسی و سه کلاس آخر برای کنترل رشته در محیط بلادرنگ در نظر گرفته شده‌اند:
\begin{center}
\eng{A, B, C, ..., Z, del, nothing, space}
\end{center}

هر کلاس دقیقاً ۳,۰۰۰ تصویر دارد؛ بنابراین تعداد کل تصاویر آموزشی برابر است با:
\[
29 \times 3000 = 87000
\]
در پوشه آزمون اولیه دیتاست تنها ۲۸ تصویر، هر کلاس یک تصویر، وجود داشته و کلاس \eng{del} غایب بوده است. این مجموعه آزمون کوچک برای ارزیابی آماری مدل کافی نبود؛ ازاین‌رو برای فاز آموزش یک تقسیم‌بندی جدید از تصاویر اصلی ایجاد شد.

\begin{figure}[H]
\centering
\safeincludegraphics[width=0.96\textwidth]{assets/class_dist.png}
\caption{توزیع کلاس‌ها در پوشه آموزشی}
\label{fig:train-dist}
\end{figure}

\section{تصویرسازی نمونه‌ها}
شکل~\ref{fig:samples} یک تصویر تصادفی از هر یک از ۲۹ کلاس را نشان می‌دهد. تصاویر شامل دست، مچ و بخشی از ساعد هستند و پس‌زمینه ثابت و قابل‌توجهی نیز در آن‌ها دیده می‌شود. این ویژگی در زمان استقرار مهم است؛ زیرا یک شبکه کانولوشنی ممکن است علاوه بر شکل دست، از نور، زاویه و الگوی پس‌زمینه نیز سرنخ بگیرد.

\begin{figure}[H]
\centering
\safeincludegraphics[width=0.98\textwidth]{assets/sample_grid.png}
\caption{نمونه‌ای از هر کلاس مجموعه آموزشی؛ تمام تصاویر \eng{200×200 RGB}}
\label{fig:samples}
\end{figure}

\section{ویژگی‌های تصاویر}
تحلیل تمام \eng{87000} تصویر نتایج زیر را نشان داد:
\begin{table}[H]
\centering
\caption{خلاصه ویژگی‌های تصویری داده}
\begin{tabularx}{0.84\textwidth}{R{5cm}Y}
\toprule
\textbf{ویژگی} & \textbf{نتیجه} \\
\midrule
حالت رنگ & همه تصاویر \eng{RGB} \\
عرض & همگی \eng{200} پیکسل \\
ارتفاع & همگی \eng{200} پیکسل \\
نسبت تصویر & ثابت و برابر \eng{1.00} \\
بازه ذخیره‌سازی پیکسل & اعداد صحیح \eng{0--255} \\
توازن کلاس‌ها & انحراف معیار تعداد نمونه‌ها برابر صفر \\
\bottomrule
\end{tabularx}
\end{table}

همان‌طور که در جدول بالا مشاهده می‌شود، تمام تصاویر موجود در مجموعه داده از ابعاد یکسان و حالت رنگی \eng{RGB} برخوردار هستند که این امر فرآیند پیش‌پردازش را به‌طور قابل‌توجهی ساده می‌سازد. همچنین داده‌ها از توازن کامل برخوردار بوده و هر یک از \eng{29} کلاس دقیقاً شامل \eng{3000} تصویر می‌باشند. این توازن کامل در کنار یکنواختی ابعاد تصاویر، مجموعه داده را به گزینه‌ای ایده‌آل برای آموزش مدل‌های یادگیری عمیق تبدیل کرده است و نیاز به تکنیک‌های متعادل‌سازی یا تغییر اندازه را به‌حداقل می‌رساند.

با وجود ثابت‌بودن اندازه داده، در هر دو مسیر آموزش و استنتاج عمل تغییر اندازه به \eng{200×200} تکرار شده است تا ورودی‌های خارج از دیتاست نیز به شکل مورد انتظار مدل تبدیل شوند.


\section{تقسیم‌بندی داده}
از کل \eng{87000} تصویر، تقسیم زیر استفاده شده است:

\begin{table}[H]
\centering
\caption{تقسیم‌بندی نهایی داده برای مدل‌سازی}
\begin{tabular}{lrr}
\toprule
\textbf{بخش} & \textbf{تعداد تصویر} & \textbf{نسبت} \\
\midrule
\eng{Train} & \eng{65250} & \eng{75\%} \\
\eng{Validation} & \eng{13050} & \eng{15\%} \\
\eng{Test} & \eng{8700} & \eng{10\%} \\
\midrule
جمع & \eng{87000} & \eng{100\%} \\
\bottomrule
\end{tabular}
\end{table}

این تقسیم‌بندی بر اساس رویه‌های استاندارد در پروژه‌های یادگیری ماشین انتخاب شده است. نسبت \eng{75} درصد برای داده‌های آموزش، حجم مناسبی از داده‌ها را برای یادگیری الگوها در اختیار مدل قرار می‌دهد. استفاده از \eng{15} درصد داده‌ها برای اعتبارسنجی، امکان پایش عملکرد مدل در حین آموزش و تنظیم فراپارامترها را بدون دخالت در فرآیند یادگیری فراهم می‌کند. همچنین در نظر گرفتن \eng{10} درصد داده‌ها به‌عنوان مجموعه آزمون، ارزیابی نهایی و بی‌طرفانه از عملکرد مدل روی داده‌های کاملاً دیده‌نشده را امکان‌پذیر می‌سازد. این نسبت‌ها به‌خوبی تعادل بین آموزش کافی، اعتبارسنجی دقیق و ارزیابی قابل‌اطمینان را برقرار می‌کنند.

\begin{figure}[H]
\centering
\safeincludegraphics[width=0.62\textwidth]{assets/dataset_split.png}
\caption{نسبت بخش‌های آموزش، اعتبارسنجی و آزمون}
\end{figure}

\section{نرمال‌سازی داده‌ها}

پس از بررسی ساختار داده‌ها، مرحله نرمال‌سازی به‌عنوان یکی از مهم‌ترین مراحل پیش‌پردازش انجام شد. هدف از این مرحله، نگاشت مقادیر پیکسل‌ها به توزیعی با میانگین صفر و واریانس واحد است که همگرایی مدل را تسریع کرده و از افت عملکرد ناشی از تغییرات شدید مقادیر ورودی جلوگیری می‌کند.

برای هر کانال رنگی، میانگین و انحراف معیار روی تمام تصاویر مجموعه آموزش محاسبه شد:
\begin{align*}
\boldsymbol{\mu} &= (0.5188341,\;0.4989899,\;0.5144102),\\
\boldsymbol{\sigma} &= (0.2045820,\;0.2336411,\;0.2412035).
\end{align*}

پس از تبدیل پیکسل‌ها به بازه \eng{[0,1]} توسط تبدیل \code{ToTensor()}، هر کانال به‌صورت جداگانه با استفاده از رابطه زیر استاندارد شد:
\[
x'_{c} = \frac{x_c - \mu_c}{\sigma_c}
\]
که در آن \(c \in \{R, G, B\}\) نشان‌دهنده کانال رنگی است.

آمار محاسبه‌شده در فایل \code{dataset\_stats.pt} ذخیره شده و در فرآیند استنتاج بلادرنگ نیز از همین مقادیر استفاده می‌شود. این همسانی در نرمال‌سازی، از تغییر توزیع داده‌های ورودی میان فاز آموزش و فاز اجرا جلوگیری کرده و یکی از عوامل کلیدی در حفظ دقت مدل در محیط عملیاتی است.

لازم به ذکر است که استفاده از آمار اختصاصی مجموعه‌داده، نسبت به استفاده از میانگین و انحراف معیار پیش‌فرض \eng{ImageNet}، عملکرد بهتری را در پی داشته است. دلیل این بهبود، نزدیکی بیشتر این آمار به توزیع واقعی داده‌های زبان اشاره است که از نظر رنگ‌بندی، بافت و نورپردازی با تصاویر طبیعی تفاوت دارد.
% ==================== Chapter 3 ====================
\chapter{فاز دوم: طراحی و آموزش شبکه عصبی پیچشی}

\section{معماری مدل}
مدل \eng{SignLanguageCNN} چهار بلوک کانولوشن و بیشینه‌گیری، یک لایه پنهان تمام‌متصل و یک لایه خروجی ۲۹کلاسه دارد. در تمام لایه‌های کانولوشن از هسته \eng{3×3}، \eng{padding=1} و تابع فعال‌ساز \eng{ReLU} استفاده شده است. هر بلوک با \eng{MaxPool 2×2} ابعاد مکانی را نصف می‌کند.

\begin{figure}[H]
\centering
\begin{tikzpicture}[
  node distance=5mm and 4mm,
  block/.style={draw=Primary,thick,rounded corners=2mm,fill=SoftBlue,minimum height=10mm,minimum width=20mm,align=center,font=\latinfont\scriptsize},
  pool/.style={draw=Secondary,thick,rounded corners=2mm,fill=SoftGreen,minimum height=10mm,minimum width=17mm,align=center,font=\latinfont\scriptsize},
  dense/.style={draw=Accent,thick,rounded corners=2mm,fill=SoftPink,minimum height=10mm,minimum width=20mm,align=center,font=\latinfont\scriptsize},
  arrow/.style={-{Stealth[length=2mm]},thick,DarkGray}
]
\begin{latin}
\node[block] (in) {Input\\$3\times200\times200$};
\node[block,right=of in] (c1) {Conv1 + ReLU\\32 channels};
\node[pool,right=of c1] (p1) {\lr{Pool}\\$32\times100\times100$};
\node[block,right=of p1] (c2) {Conv2 + ReLU\\64 channels};
\node[pool,right=of c2] (p2) {\lr{Pool}\\$64\times50\times50$};
\node[block,below=9mm of p2] (c3) {Conv3 + ReLU\\128 channels};
\node[pool,left=of c3] (p3) {\lr{Pool}\\$128\times25\times25$};
\node[block,left=of p3] (c4) {Conv4 + ReLU\\256 channels};
\node[pool,left=of c4] (p4) {\lr{Pool}\\$256\times12\times12$};
\node[dense,left=of p4] (fc1) {Flatten + FC\\$36864\rightarrow512$};
\node[dense,left=of fc1] (drop) {Dropout\\$p=0.5$};
\node[dense,left=of drop] (out) {FC Output\\29 logits};
\draw[arrow] (in)--(c1); \draw[arrow] (c1)--(p1); \draw[arrow] (p1)--(c2); \draw[arrow] (c2)--(p2);
\draw[arrow] (p2)--(c3); \draw[arrow] (c3)--(p3); \draw[arrow] (p3)--(c4); \draw[arrow] (c4)--(p4);
\draw[arrow] (p4)--(fc1); \draw[arrow] (fc1)--(drop); \draw[arrow] (drop)--(out);
\end{latin}
\end{tikzpicture}
\caption{معماری شبکه \eng{SignLanguageCNN}}
\end{figure}

\begin{table}[H]
\centering
\caption{جزئیات لایه‌ها و تعداد پارامترهای قابل آموزش}
\begin{tabularx}{\textwidth}{l C{2.8cm} C{2.8cm} r Y}
\toprule
\textbf{لایه} & \textbf{ورودی} & \textbf{خروجی پس از \eng{Pool}} & \textbf{پارامتر} & \textbf{توضیح} \\
\midrule
\eng{Conv1} & \eng{3×200×200} & \eng{32×100×100} & \eng{896} & استخراج لبه‌ها و الگوهای ابتدایی \\
\eng{Conv2} & \eng{32×100×100} & \eng{64×50×50} & \eng{18,496} & ترکیب ویژگی‌های سطح پایین \\
\eng{Conv3} & \eng{64×50×50} & \eng{128×25×25} & \eng{73,856} & یادگیری الگوهای پیچیده‌تر شکل دست \\
\eng{Conv4} & \eng{128×25×25} & \eng{256×12×12} & \eng{295,168} & تولید نمایش فشرده سطح بالا \\
\eng{FC1} & \eng{36,864} & \eng{512} & \eng{18,874,880} & تبدیل ویژگی‌های مکانی به نمایش طبقه‌بندی \\
\eng{FC2} & \eng{512} & \eng{29} & \eng{14,877} & تولید لاجیت هر کلاس \\
\midrule
\multicolumn{3}{r}{\textbf{مجموع}} & \eng{19,278,173} & \\
\bottomrule
\end{tabularx}
\end{table}
\subsection{دلایل انتخاب اجزای معماری}

\subsubsection{تعداد لایه‌های کانولوشن}

تعداد \eng{4} لایه کانولوشن بر اساس ویژگی‌های داده و آزمایش‌های اولیه انتخاب شده است. تصاویر حروف اشاره شامل الگوهای محدودی هستند و \eng{4} لایه برای استخراج ویژگی‌های ساده تا پیچیده کافی است. همچنین با توجه به ابعاد \eng{200×200} پیکسل، پس از \eng{4} بار بیشینه‌گیری، ابعاد به \eng{12×12} کاهش می‌یابد که حفظ اطلاعات مکانی را تضمین می‌کند. آزمایش‌های اولیه با \eng{3} و \eng{5} لایه نشان داد که \eng{3} لایه دقت کافی ندارد و \eng{5} لایه نیز بدون بهبود چشمگیر، تعداد پارامترها و خطر بیش‌برازش را افزایش می‌دهد. بنابراین \eng{4} لایه به‌عنوان معماری نهایی انتخاب شده است.

\subsubsection{هسته \eng{3×3} و \eng{padding=1}}

هسته \eng{3×3} کوچک‌ترین هسته مؤثر برای تشخیص الگوهای محلی مانند لبه‌ها و بافت‌ها است و تعداد پارامتر کمی دارد. \eng{padding=1} نیز باعث می‌شود ابعاد تصویر در طول کانولوشن حفظ شود و اطلاعات لبه‌ها از دست نرود. این ترکیب، تعادل مناسبی بین قدرت تشخیص و حجم محاسبات ایجاد می‌کند.

\subsubsection{افزایش تدریجی تعداد کانال‌ها}

تعداد کانال‌ها از \eng{32} به \eng{256} افزایش یافته است. لایه‌های اولیه برای تشخیص الگوهای ساده مانند لبه‌ها به کانال کمتری نیاز دارند، درحالی‌که لایه‌های عمیق‌تر برای ترکیب ویژگی‌ها و تشخیص الگوهای پیچیده‌تر مانند شکل انگشتان، به کانال بیشتری نیاز دارند. این افزایش تدریجی، تعادل بین قدرت تشخیص و حجم محاسبات را حفظ می‌کند.

\subsubsection{لایه \eng{MaxPool 2×2}}

لایه \eng{MaxPool 2×2} ابعاد مکانی را نصف می‌کند و باعث کاهش حجم داده، افزایش مقاومت در برابر جابه‌جایی جزئی دست و انتخاب ویژگی‌های مهم می‌شود. این کار همچنین تعداد پارامترهای لایه‌های بعدی را کاهش داده و سرعت پردازش را افزایش می‌دهد.

\subsubsection{تابع فعال‌ساز \eng{ReLU}}

تابع \eng{ReLU} به دلیل سادگی، سرعت بالا و کاهش مشکل محو شدن گرادیان در شبکه‌های عمیق انتخاب شده است. این تابع از نظر محاسباتی بسیار کارآمد است و در برابر اشباع شدن نورون‌ها مقاومت خوبی نشان می‌دهد.

\subsubsection{لایه \eng{FC1} با \eng{512} نورون}

\eng{512} نورون در لایه پنهان تمام‌متصل، ظرفیت کافی برای یادگیری ترکیبات پیچیده ویژگی‌ها را فراهم می‌کند. این تعداد نه‌چندان کم است که قدرت طبقه‌بندی را محدود کند و نه‌چندان زیاد که باعث بیش‌برازش شود. با توجه به تعداد کلاس‌ها (\eng{29} کلاس)، این تعداد نورون انتخاب مناسبی است.

\subsubsection{\eng{Dropout} با نرخ \eng{0.5}}

با توجه به تعداد بالای پارامترها در لایه \eng{FC1} (حدود \eng{19} میلیون)، از چندین روش منظم‌سازی برای جلوگیری از بیش‌برازش استفاده شده است. \eng{Dropout} با نرخ \eng{0.5} به‌عنوان یکی از این روش‌ها، به‌خوبی خطر بیش‌برازش را کاهش می‌دهد و در عین حال اجازه می‌دهد مدل اطلاعات مفید را یاد بگیرد. همچنین استفاده از \eng{Data Augmentation} و کاهش نرخ یادگیری در زمان مناسب، از دیگر اقداماتی است که به پایداری و تعمیم‌پذیری مدل کمک کرده است.

\subsubsection{لایه خروجی \eng{29} نورون}

تعداد نورون‌های لایه خروجی دقیقاً برابر با تعداد کلاس‌های مسئله (\eng{29} کلاس: \eng{A-Z, del, nothing, space}) است. هر نورون نشان‌دهنده احتمال تعلق تصویر ورودی به یکی از کلاس‌ها است و خروجی نهایی با تابع \eng{Softmax} به توزیع احتمال تبدیل می‌شود.


\section{افزایش داده و کنترل بیش‌برازش}

به منظور افزایش قدرت تعمیم‌دهی مدل و جلوگیری از بیش‌برازش، از مجموعه‌ای از تبدیلات تصادفی در مرحله آموزش استفاده شده است. این تبدیلات شامل تغییر اندازه به \eng{200×200}، چرخش تصادفی تا \lr{$\pm 10^\circ$}، وارونگی افقی تصادفی و تغییر جزئی روشنایی و کنتراست با شدت \eng{0.15} است. 

لازم به ذکر است که در نسخه نهایی پیاده‌سازی، به دلیل تأثیر منفی وارونگی افقی روی برخی کلاس‌های نامتقارن مانند \eng{Z}، \eng{J}، \eng{R} و \eng{P}، این تبدیل از فرآیند آموزش حذف شده است. با حذف این تبدیل، دقت مدل روی کلاس‌های مذکور بهبود قابل‌توجهی نشان داد. بنابراین تبدیلات نهایی اعمال‌شده در مرحله آموزش به شرح زیر است:

\begin{lstlisting}[language=Python]
transform_train = transforms.Compose([
    transforms.Resize((200, 200)),
    transforms.RandomRotation(10),
    transforms.ColorJitter(brightness=0.15, contrast=0.15),
    transforms.ToTensor(),
    transforms.Normalize(mean=mean, std=std)
])
\end{lstlisting}

پس از اعمال این تبدیلات، تصاویر به تنسور تبدیل شده و با استفاده از آمار محاسبه‌شده از داده‌های آموزش (میانگین و انحراف معیار هر کانال) نرمال‌سازی می‌شوند. برای داده‌های اعتبارسنجی و آزمون، تنها تغییر اندازه، تبدیل به تنسور و نرمال‌سازی انجام می‌شود تا ارزیابی روی داده‌های بدون اغتشاش صورت گیرد:

\begin{lstlisting}[language=Python]
transform_val = transforms.Compose([
    transforms.Resize((200, 200)),
    transforms.ToTensor(),
    transforms.Normalize(mean=mean, std=std)
])
\end{lstlisting}

\subsection{تکنیک‌های استفاده‌شده برای بهبود تعمیم}
در معماری نهایی، سه تکنیک اصلی برای کنترل بیش‌برازش و بهبود تعمیم‌دهی پیاده‌سازی شده است:

\begin{enumerate}
    \item \textbf{\eng{Data Augmentation}:} با استفاده از تبدیلات تصادفی (چرخش و تغییرات رنگ)، تنوع داده‌های آموزش به‌طور مصنوعی افزایش یافته است. این تبدیلات با دقت انتخاب شده‌اند تا ضمن افزایش تنوع، از ایجاد تغییرات غیرطبیعی که منجر به کاهش دقت می‌شوند، جلوگیری کنند.
    \item \textbf{\eng{Dropout}:} لایه \eng{Dropout} با نرخ \eng{0.5} پس از لایه \eng{FC1} قرار گرفته است. با توجه به تعداد بالای پارامترها در این لایه (حدود \eng{19} میلیون)، این تکنیک به‌خوبی وابستگی نورون‌ها به یکدیگر را کاهش داده و از بیش‌برازش جلوگیری می‌کند.
    \item \textbf{\eng{ReduceLROnPlateau}:} نرخ یادگیری در صورت توقف بهبود زیان اعتبارسنجی به‌مدت \eng{3} دوره، به‌صورت خودکار با ضریب \eng{0.5} کاهش می‌یابد. این کار به مدل کمک می‌کند تا از نقاط بهینه محلی عبور کرده و به نقاط بهینه سراسری نزدیک‌تر شود.
\end{enumerate}

\section{پارامترهای آموزش}

پارامترهای استفاده‌شده در فرآیند آموزش در جدول زیر خلاصه شده است:

\begin{table}[H]
\centering
\caption{ابرپارامترهای آموزش}
\begin{tabularx}{0.86\textwidth}{R{5.2cm}Y}
\toprule
\textbf{پارامتر} & \textbf{مقدار} \\
\midrule
تعداد دوره & \eng{10} \\
اندازه دسته & \eng{64} \\
نرخ یادگیری اولیه & \eng{0.001} \\
بهینه‌ساز & \eng{Adam} \\
تابع زیان & \eng{CrossEntropyLoss} \\
زمان‌بند نرخ یادگیری & \eng{ReduceLROnPlateau} \\
تنظیم زمان‌بند & \eng{patience=3, factor=0.5} \\
معیار ذخیره بهترین مدل & بیشترین دقت اعتبارسنجی \\
فایل وزن‌ها & \code{models/final_model.pth} \\
\bottomrule
\end{tabularx}
\end{table}

نرخ یادگیری اولیه ۰.۰۰۱ با بهینه‌ساز \eng{Adam} انتخاب شده است. تابع زیان \eng{CrossEntropyLoss} برای مسائل طبقه‌بندی چندکلاسه مناسب است و خروجی مدل (لاجیت‌ها) را با برچسب‌های یک‌داغ مقایسه می‌کند. زمان‌بند \eng{ReduceLROnPlateau} در صورت عدم بهبود زیان اعتبارسنجی به مدت ۳ دوره، نرخ یادگیری را به اندازه ضریب ۰.۵ کاهش می‌دهد. این کار به مدل کمک می‌کند تا در نقاط بهینه محلی گیر نکند و به نقاط بهینه سراسری نزدیک‌تر شود. اندازه دسته ۶۴ با توجه به حافظه \eng{GPU} و تعادل بین سرعت و پایداری گرادیان انتخاب شده است. بهترین مدل بر اساس بیشترین دقت اعتبارسنجی انتخاب و در فایل \code{models/final_model.pth} ذخیره شده است.

\subsection{انتخاب بهینه‌ساز}

بهینه‌ساز \lr{Adam} به دلیل مزایای زیر برای آموزش مدل انتخاب شده است:

\begin{itemize}
    \item \textbf{سرعت همگرایی بالا:} \eng{Adam} با ترکیب \eng{Momentum} و \eng{Adaptive Learning Rate}، همگرایی سریع‌تری نسبت به \eng{SGD} دارد و با تعداد دوره محدود (۱۰ \eng{Epoch}) به دقت مناسب می‌رسد.
    \item \textbf{تنظیم خودکار نرخ یادگیری:} این بهینه‌ساز برای هر پارامتر، نرخ یادگیری جداگانه تنظیم می‌کند و نیاز به تنظیم دقیق نرخ یادگیری اولیه را کاهش می‌دهد.
    \item \textbf{مقاومت در برابر نویز:} داده‌های تصویری شامل نویز هستند و \eng{Adam} با استفاده از میانگین متحرک گرادیان‌ها، این نویز را کاهش می‌دهد.
    \item \textbf{مناسب برای داده‌های بزرگ:} با توجه به حجم بالای داده‌های آموزش (۸۷,۰۰۰ تصویر)، \eng{Adam} عملکرد پایدار و مناسبی ارائه می‌دهد.
\end{itemize}

اگرچه \eng{SGD} نیز گزینه‌ای رایج است، اما به دلیل همگرایی کندتر و نیاز به تنظیم دقیق‌تر نرخ یادگیری، برای این پروژه با محدودیت تعداد دوره، انتخاب مناسبی نبود. تجربه عملی نیز نشان داده است که \eng{Adam} در اکثر پروژه‌های بینایی ماشین به نتایج بهتری می‌رسد.

\section{روند آموزش}

\begin{figure}[H]
\centering
\safeincludegraphics[width=0.88\textwidth]{assets/training_accuracy.png}
\caption{دقت آموزش و اعتبارسنجی}
\label{fig:acc}
\end{figure}

\begin{table}[H]
\centering
\caption{دقت ثبت‌شده در هر دوره}
\begin{tabular}{crr|crr}
\toprule
\textbf{دوره} & \textbf{آموزش} & \textbf{اعتبارسنجی} & \textbf{دوره} & \textbf{آموزش} & \textbf{اعتبارسنجی} \\
\midrule
1 & \eng{64.38\%} & \eng{96.00\%} & 6 & \eng{97.53\%} & \eng{99.43\%} \\
2 & \eng{91.56\%} & \eng{98.77\%} & 7 & \eng{97.86\%} & \eng{99.57\%} \\
3 & \eng{95.01\%} & \eng{99.09\%} & 8 & \eng{98.22\%} & \eng{99.71\%} \\
4 & \eng{96.29\%} & \eng{99.48\%} & 9 & \eng{98.26\%} & \eng{99.87\%} \\
5 & \eng{97.02\%} & \eng{99.56\%} & 10 & \eng{98.34\%} & \eng{99.91\%} \\
\bottomrule
\end{tabular}
\end{table}
در ابتدا، فرآیند آموزش با \eng{30} دوره برنامه‌ریزی شده بود. اما با بررسی روند تغییرات دقت و زیان در دوره‌های ابتدایی، مشخص شد که مدل به‌سرعت به دقت بالایی روی داده‌های اعتبارسنجی دست می‌یابد و از دوره دوم به بعد، دقت اعتبارسنجی به بیش از \eng{98} درصد می‌رسد. همچنین از دوره سوم به بعد، دقت اعتبارسنجی همواره بالای \eng{99} درصد بوده و در دوره نهم و دهم به ترتیب به \eng{99.87} و \eng{99.91} درصد رسیده است. دقت آموزش نیز به‌تدریج از \eng{64.38} درصد در دوره اول به \eng{98.34} درصد در دوره دهم افزایش یافته است.

با توجه به رشد سریع دقت در دوره‌های ابتدایی و تثبیت آن در محدوده \eng{99} درصد، تعداد دوره‌ها از \eng{30} به \eng{10} کاهش یافت تا از اتلاف وقت و افزایش ریسک بیش‌برازش جلوگیری شود. انتخاب \eng{10} دوره به‌عنوان تعداد نهایی دوره‌ها، تعادل مناسبی بین دستیابی به دقت بالا و جلوگیری از بیش‌برازش ایجاد می‌کند.

همان‌طور که در جدول و نمودار مشاهده می‌شود، بهترین دقت اعتبارسنجی در دوره دهم با مقدار \eng{99.91} درصد ثبت شده است. بنابراین وزن‌های مربوط به این دوره به‌عنوان بهترین مدل ذخیره شده و در ادامه مراحل استنتاج و ارزیابی نهایی مورد استفاده قرار گرفته است. عملکرد مدل در دوره‌های پایانی نشان‌دهنده پایداری مناسب مدل و تأیید انتخاب \eng{10} دوره به‌عنوان تعداد مناسب است.

\section{ارزیابی نهایی}
نتایج ثبت‌شده روی مجموعه آزمون \eng{8700} تصویری عبارت‌اند از:
\begin{table}[H]
\centering
\caption{نتایج نهایی مدل}
\begin{tabularx}{0.78\textwidth}{R{5cm}Y}
\toprule
\textbf{معیار} & \textbf{مقدار} \\
\midrule
پیش‌بینی صحیح & \eng{8,695 / 8,700} \\
دقت آزمون & \eng{99.94\%} \\
میانگین اطمینان & \eng{99.89\%} \\
کمینه اطمینان مشاهده‌شده & \eng{40.62\%} \\
\bottomrule
\end{tabularx}
\end{table}

دقت \eng{99.94} درصد روی مجموعه آزمون نشان‌دهنده عملکرد بسیار خوب مدل در تشخیص حروف اشاره است. این نتیجه حاکی از آن است که مدل الگوهای موجود در داده‌های آموزشی را به‌خوبی یاد گرفته و توانسته است روی داده‌های دیده‌نشده نیز تعمیم‌پذیری بالایی داشته باشد. میانگین اطمینان \eng{99.89} درصد نیز بیانگر آن است که مدل در اکثر پیش‌بینی‌ها با اطمینان بالا تصمیم‌گیری می‌کند. کمینه اطمینان مشاهده‌شده (\eng{40.62} درصد) مربوط به تعداد بسیار کمی از تصاویر با کیفیت پایین یا نویز بالا است که در ارزیابی کلی تأثیر قابل‌توجهی ندارد.

\begin{figure}[H]
\centering
\safeincludegraphics[width=0.9\textwidth]{assets/confusion_matrix.png}
\caption{ماتریس درهم‌ریختگی روی مجموعه آزمون}
\label{fig:confusion}
\end{figure}

ماتریس درهم‌ریختگی نشان می‌دهد که مدل در اکثر کلاس‌ها عملکرد عالی داشته و تعداد خطاها بسیار محدود است. بیشترین خطاها مربوط به کلاس‌های مشابه زیر است:
\begin{table}[H]
\centering
\caption{پنج جفت کلاس با بیشترین خطا}
\begin{tabular}{c|c|c}
\toprule
\textbf{کلاس واقعی} & \textbf{پیش‌بینی اشتباه} & \textbf{تعداد خطا} \\
\midrule
\eng{N} & \eng{M} & ۲ \\
\eng{V} & \eng{W} & ۱ \\
\eng{del} & \eng{B} & ۱ \\
\eng{del} & \eng{C} & ۱ \\
\bottomrule
\end{tabular}
\end{table}
این خطاها عمدتاً بین کلاس‌هایی رخ داده است که از نظر شکل دست و موقعیت انگشتان شباهت زیادی به یکدیگر دارند. برای مثال، حروف \eng{N} و \eng{M} در زبان اشاره تنها در تعداد انگشتان باز شده تفاوت دارند و این شباهت بصری باعث سردرگمی مدل شده است. همچنین حروف \eng{V} و \eng{W} نیز از نظر زاویه انگشتان بسیار نزدیک به هم هستند. مجموع خطاها \eng{5} مورد از \eng{8700} تصویر (معادل \eng{0.06} درصد) است که نشان‌دهنده عملکرد بسیار دقیق مدل است.

\section{بهبود مدل با \eng{Weighted Loss} (تنظیم دقیق)}

پس از ارزیابی مدل روی داده‌های تست، عملکرد بسیار مطلوبی با دقت \eng{99.34}\% به‌دست آمد. با این حال، پس از اجرای مدل در محیط واقعی (فاز سوم) و تست روی تصاویر وب‌کم، مشاهده شد که مدل روی برخی کلاس‌ها عملکرد ضعیف‌تری دارد و آنها را با کلاس‌های مشابه اشتباه می‌گیرد. این کلاس‌ها عمدتاً شامل حروفی هستند که از نظر شکل دست و حالت انگشتان شباهت زیادی به یکدیگر دارند. برای رفع این مشکل، از تکنیک \eng{Weighted Loss} استفاده شد تا مدل توجه بیشتری به کلاس‌های ضعیف داشته باشد و بتواند تفاوت‌های ظریف بین آنها را بهتر یاد بگیرد.

\subsection{مرحله ۱: وزن‌دهی به کلاس‌های \eng{A}، \eng{M} و \eng{N}}

کلاس‌های \eng{A}، \eng{M} و \eng{N} از جمله کلاس‌هایی بودند که در تشخیص \eng{Real-time} بیشترین خطا را داشتند. هر سه این حروف در زبان اشاره به‌صورت مشت با تغییرات جزئی در موقعیت شست و انگشتان اجرا می‌شوند که تشخیص آنها را دشوار می‌کند. به همین دلیل، وزن این کلاس‌ها در تابع هزینه به مقدار \eng{2.5} افزایش یافت. سپس مدل با این وزن‌های جدید به مدت \eng{5} دوره (\eng{Epoch}) با نرخ یادگیری \eng{0.0001} تنظیم دقیق شد. پس از این مرحله، دقت مدل روی این سه کلاس به‌طور محسوسی بهبود یافت و خطای تشخیص آنها کاهش پیدا کرد.

\subsection{مرحله ۲: وزن‌دهی به کلاس‌های \eng{A}، \eng{S} و \eng{T}}

در مرحله بعد، کلاس‌های \eng{A}، \eng{S} و \eng{T} مورد بررسی قرار گرفتند. این کلاس‌ها نیز به دلیل شباهت در حالت مشت، در محیط \eng{Real-time} با چالش مواجه بودند. به‌ویژه کلاس \eng{S} که تشخیص آن در شرایط نوری مختلف با دشواری همراه بود. بنابراین، وزن این کلاس‌ها نیز به مقدار \eng{2.5} تنظیم شد. این مرحله با \eng{3} دوره تنظیم دقیق انجام شد تا مدل تفاوت‌های ظریف بین این کلاس‌ها را بهتر تشخیص دهد. نتایج نشان داد که دقت تشخیص این کلاس‌ها به‌ویژه \eng{S} و \eng{T} در تصاویر \eng{Real-time} بهبود یافته است.

\subsection{مرحله ۳: یکسان‌سازی وزن‌ها}

پس از بهبود کلاس‌های هدف، برای جلوگیری از بیش‌برازش و حفظ تعادل مدل، وزن‌های همه کلاس‌ها به مقدار \eng{1} برگردانده شد. این کار باعث شد مدل بدون اولویت‌دهی خاص، روی همه کلاس‌ها به‌صورت یکسان تمرکز کند و عملکرد کلی آن پایدار بماند. در نهایت، مدل با وزن‌های یکسان به مدت \eng{3} دوره دیگر تنظیم دقیق شد تا وزن‌های نهایی تثبیت شوند و مدل برای استفاده در محیط \eng{Real-time} آماده گردد.

\subsection{نتایج حاصل از تنظیم دقیق}
جدول زیر خلاصه‌ای از مراحل انجام‌شده و نتایج حاصل را نشان می‌دهد:
\begin{table}[H]
\centering
\caption{خلاصه مراحل Fine-Tuning و نتایج}
\begin{tabularx}{0.9\textwidth}{C{2.5cm}C{3cm}C{2.5cm}C{3cm}}
\toprule
\textbf{مرحله} & \textbf{کلاس‌های هدف} & \textbf{تعداد دوره} & \textbf{نتیجه} \\
\midrule
مرحله ۱ & \eng{A, M, N} & \eng{5} & بهبود تشخیص این کلاس‌ها \\
مرحله ۲ & \eng{A, S, T} & \eng{3} & بهبود تشخیص \eng{S} و \eng{T} در Real-time \\
مرحله ۳ & همه کلاس‌ها (وزن یکسان) & \eng{3} & تثبیت عملکرد و جلوگیری از Overfitting \\
\bottomrule
\end{tabularx}
\end{table}

در نهایت، با استفاده از این رویکرد گام‌به‌گام، مدل توانست بدون کاهش دقت روی کلاس‌های قوی، عملکرد خود را روی کلاس‌های ضعیف به‌طور قابل‌توجهی بهبود بخشد و برای استفاده در محیط واقعی آماده شود. وزن‌های نهایی مدل در فایل \code{models/final\_model.pth} ذخیره شده است.
% ==================== Chapter 4 ====================
\chapter{فاز سوم: سامانه بلادرنگ ترجمه به متن و گفتار}

\section{معماری کلان}
کد فاز سوم به‌صورت یک بسته ماژولار با چهار حوزه اصلی طراحی شده است: بینایی ماشین، استنتاج مدل، رمزگشایی زمانی و گفتار. جریان کامل داده در شکل~\ref{fig:pipeline} آمده است.

\begin{figure}[H]
\centering
\resizebox{0.98\textwidth}{!}{%
\begin{tikzpicture}[
  box/.style={draw=Primary,thick,rounded corners=2mm,fill=SoftBlue,minimum width=25mm,minimum height=13mm,align=center},
  model/.style={draw=Accent,thick,rounded corners=2mm,fill=SoftPink,minimum width=26mm,minimum height=13mm,align=center},
  logic/.style={draw=Secondary,thick,rounded corners=2mm,fill=SoftGreen,minimum width=29mm,minimum height=13mm,align=center},
  arrow/.style={-{Stealth[length=3mm]},very thick,DarkGray},
  node distance=8mm
]
\begin{latin}
\node[box] (cam) {Webcam\\OpenCV};
\node[box,right=of cam] (roi) {Mirror + Square ROI};
\node[box,right=of roi] (pre) {Resize, RGB, Normalize\\NCHW batch};
\node[model,right=of pre] (cnn) {PyTorch CNN\\29 logits};
\node[model,right=of cnn] (soft) {Softmax\\Prediction};
\node[logic,below=13mm of soft] (temp) {Temporal Decoder\\2 s hold};
\node[logic,left=of temp] (handler) {Label Handler\\A--Z / SPACE / DEL};
\node[logic,left=of handler] (buffer) {Text Buffer};
\node[logic,left=of buffer] (final) {Sequence Finalizer\\5 s NOTHING};
\node[box,left=of final] (tts) {pyttsx3\\Speech};
\draw[arrow] (cam)--(roi); \draw[arrow] (roi)--(pre); \draw[arrow] (pre)--(cnn); \draw[arrow] (cnn)--(soft);
\draw[arrow] (soft)--(temp); \draw[arrow] (temp)--(handler); \draw[arrow] (handler)--(buffer); \draw[arrow] (buffer)--(final); \draw[arrow] (final)--(tts);
\end{latin}
\end{tikzpicture}}
\caption{زنجیره پردازش بلادرنگ از وب‌کم تا گفتار}
\label{fig:pipeline}
\end{figure}

\section{ساختار ماژول‌ها}
\begin{latin}
\begin{lstlisting}[style=pythonstyle,language={},numbers=none]
src/sign_translator/
|-- app.py                 # orchestration and main loop
|-- config.py              # camera/model/decoder/TTS parameters
|-- labels.py              # canonical label order and commands
|-- vision/
|   |-- camera.py          # safe camera lifecycle
|   |-- roi.py             # square region of interest
|   `-- overlay.py         # prediction and progress UI
|-- inference/
|   |-- preprocessing.py   # BGR -> normalized NCHW
|   |-- torch_model.py     # SignLanguageCNN architecture
|   |-- torch_predictor.py # weight loading and inference
|   |-- output_decoder.py  # logits/probabilities validation
|   `-- factory.py         # mock or torch backend
|-- decoding/
|   |-- temporal_decoder.py
|   |-- sequence_finalizer.py
|   |-- label_handler.py
|   `-- text_buffer.py
`-- speech/
    |-- factory.py
    |-- noop.py
    `-- pyttsx3_engine.py
\end{lstlisting}
\end{latin}

این جداسازی باعث می‌شود مدل، موتور گفتار و منطق رابط بدون بازنویسی حلقه اصلی قابل تعویض باشند. برای مثال، \code{TimedMockPredictor} پیش از تحویل مدل واقعی امکان تست کل سامانه را فراهم کرده و سپس \code{TorchPredictor} از طریق کارخانه جایگزین آن شده است.

\section{دریافت تصویر و ناحیه مورد توجه}
کلاس \code{Camera} دوربین را به‌صورت \eng{context manager} باز و در پایان آزاد می‌کند. رزولوشن در تنظیمات روی \eng{1280×720} قرار گرفته و تصویر برای تجربه طبیعی کاربر به‌صورت افقی آینه می‌شود.

ناحیه دست با نسبت‌های مستقل از رزولوشن تعریف شده است:
\begin{align*}
 x &= 0.55W, & y &= 0.12H,\\
 w &= 0.38W, & h &= 0.72H.
\end{align*}
سپس ضلع کوچک‌تر انتخاب و کادر دقیقاً مربعی می‌شود. این اقدام از فشردگی یا کشیدگی شکل دست هنگام تبدیل به ورودی \eng{200×200} جلوگیری می‌کند.

\section{همسان‌سازی پیش‌پردازش آموزش و اجرا}
فریم \eng{OpenCV} در قالب \eng{BGR} است، درحالی‌که داده آموزشی \eng{RGB} بوده است. مسیر استنتاج مراحل زیر را انجام می‌دهد:
\begin{enumerate}
  \item اعتبارسنجی ورودی سه‌کاناله و غیرخالی
  \item تغییر اندازه به \eng{200×200}
  \item تبدیل \eng{BGR} به \eng{RGB}
  \item تبدیل نوع به \eng{float32} و تقسیم بر \eng{255}
  \item نرمال‌سازی کانالی با میانگین و انحراف معیار آموزش
  \item تبدیل چیدمان از \eng{HWC} به \eng{CHW}
  \item افزودن بعد دسته و تولید تنسور \eng{1×3×200×200}
\end{enumerate}
انتخاب روش درون‌یابی نیز بر اساس کوچک‌سازی یا بزرگ‌سازی انجام می‌شود: \eng{INTER\_AREA} برای کاهش اندازه و \eng{INTER\_LINEAR} برای افزایش اندازه.

\section{بارگذاری مدل و استنتاج}
\code{TorchPredictor} معماری را با تعداد خروجی برابر طول فهرست برچسب‌ها می‌سازد و فایل وزن را با \code{weights_only=True} بارگذاری می‌کند. استفاده از \code{strict=True} در \code{load_state_dict} سبب می‌شود هر ناسازگاری میان معماری و وزن‌ها به‌جای تولید نتیجه اشتباه، با خطای واضح متوقف شود.

ترتیب خروجی مدل در کد به‌صورت زیر ثابت شده است:
\begin{center}
\eng{(A, ..., Z, DEL, NOTHING, SPACE)}
\end{center}
خروجی شبکه لاجیت است. \code{output_decoder.py} شکل خروجی، تعداد کلاس‌ها، مقادیر نامتناهی و تکراری‌نبودن برچسب‌ها را بررسی و سپس \eng{softmax} پایدار عددی را محاسبه می‌کند.

\section{منطق تثبیت زمانی}
پیش‌بینی خام هر فریم مستقیماً وارد متن نمی‌شود. \code{TemporalDecoder} یک ماشین حالت ساده دارد:
\begin{itemize}
  \item پیش‌بینی باید اطمینان حداقل \eng{0.80} داشته باشد
  \item یک برچسب باید به‌مدت \eng{2.0 s} بدون تغییر باقی بماند
  \item پس از ثبت، همان برچسب قفل می‌شود تا با نگه‌داشتن دست چندبار تکرار نشود
  \item نمایش \eng{NOTHING} یا افت اطمینان به‌مدت \eng{0.4 s} قفل را آزاد می‌کند
  \item برای تکرار دو حرف یکسان، مانند \eng{OO}، کاربر باید بین دو نمایش، حالت \eng{NOTHING} را نشان دهد
\end{itemize}

\begin{figure}[H]
\centering
\begin{tikzpicture}[
  state/.style={draw=Primary,thick,rounded corners=3mm,fill=SoftBlue,minimum width=34mm,minimum height=13mm,align=center},
  accept/.style={draw=Accent,thick,rounded corners=3mm,fill=SoftPink,minimum width=34mm,minimum height=13mm,align=center},
  neutral/.style={draw=Secondary,thick,rounded corners=3mm,fill=SoftGreen,minimum width=34mm,minimum height=13mm,align=center},
  arrow/.style={-{Stealth[length=2.5mm]},thick,DarkGray},node distance=20mm
]
\begin{latin}
\node[state] (wait) {Waiting};
\node[state,right=of wait] (cand) {Candidate\\confidence $\geq0.80$};
\node[accept,right=of cand] (lock) {Accepted / Locked};
\node[neutral,below=of lock] (rel) {NOTHING / low confidence\\for $0.4$ s};
\draw[arrow] (wait) -- node[above,font=\scriptsize]{new sign} (cand);
\draw[arrow] (cand) -- node[above,font=\scriptsize]{stable for 2 s} (lock);
\draw[arrow] (cand) edge[bend left=30] node[below,font=\scriptsize]{label changes} (wait);
\draw[arrow] (lock) -- (rel);
\draw[arrow] (rel) -| (wait);
\end{latin}
\end{tikzpicture}
\caption{ماشین حالت رمزگشایی زمانی}
\end{figure}

\section{تفسیر برچسب‌ها و بافر متن}
\begin{table}[H]
\centering
\caption{رفتار هر نوع برچسب پذیرفته‌شده}
\begin{tabularx}{0.84\textwidth}{C{3cm}Y}
\toprule
\textbf{برچسب} & \textbf{عملکرد} \\
\midrule
\eng{A--Z} & افزودن حرف انگلیسی بزرگ به انتهای بافر \\
\eng{SPACE} & افزودن یک فاصله؛ فاصله ابتدایی یا تکراری نادیده گرفته می‌شود \\
\eng{DEL} & حذف آخرین کاراکتر، شامل حرف یا فاصله \\
\eng{NOTHING} & عدم تغییر متن؛ استفاده به‌عنوان علامت خنثی و پایان رشته \\
\bottomrule
\end{tabularx}
\end{table}

\code{TextBuffer} به‌جای ذخیره توکن‌های کلمه، فهرست کاراکترها را نگه می‌دارد. بنابراین فرمان حذف دقیقاً آخرین کاراکتر را برمی‌گرداند و متن نهایی با حذف فاصله‌های ابتدا و انتها تولید می‌شود.

\section{تشخیص پایان عبارت و تبدیل به گفتار}
\code{SequenceFinalizer} تنها زمانی شمارنده را فعال می‌کند که:
\begin{itemize}
  \item برچسب \eng{NOTHING} با اطمینان حداقل \eng{0.80} تشخیص داده شود.
  \item بافر متن خالی نباشد.
  \item این حالت برای \eng{5.0 s} پیوسته ادامه یابد.
\end{itemize}
پس از تکمیل شمارنده، متن یک‌بار نهایی می‌شود، به موتور گفتار ارسال و سپس بافر پاک می‌شود. پرچم داخلی از چندبارخواندن متن در همان دوره طولانی \eng{NOTHING} جلوگیری می‌کند.

موتور \eng{pyttsx3} در یک نخ کارگر اجرا می‌شود و متن‌ها را از صف دریافت می‌کند. این طراحی مانع توقف حلقه پردازش ویدئو هنگام پخش صدا می‌شود. یک پیاده‌سازی \eng{NoOp} نیز برای تست و محیط‌های فاقد صوت وجود دارد.

\section{رابط تصویری}
\code{overlay.py} اطلاعات زیر را روی فریم نشان می‌دهد:
\begin{itemize}
  \item کادر سبز ناحیه دست
  \item برچسب پیش‌بینی‌شده و درصد اطمینان
  \item درصد پیشرفت نگه‌داشتن علامت
  \item وضعیت قفل، پذیرش یا آزادشدن علامت
  \item درصد پیشرفت نهایی‌سازی متن
  \item متن جاری و آخرین متن نهایی‌شده
  \item راهنمای خروج با کلیدهای \eng{Q} یا \eng{ESC}
\end{itemize}

% ==================== Chapter 5 ====================
\chapter{اعتبارسنجی نرم‌افزار و کیفیت پیاده‌سازی}

\section{مجموعه تست‌ها}
پوشه \code{tests} اجزای زیر را پوشش می‌دهد:
\begin{table}[H]
\centering
\caption{حوزه‌های تحت آزمون واحد}
\begin{tabularx}{\textwidth}{R{5cm}Y}
\toprule
\textbf{فایل تست} & \textbf{رفتارهای بررسی‌شده} \\
\midrule
\code{test_roi.py} & محاسبه کادر، برش تصویر و مربعی‌بودن ناحیه ورودی \\
\code{test_preprocessing.py} & شکل \eng{NCHW}، تبدیل رنگ، نوع داده و نرمال‌سازی کانال‌ها \\
\code{test_torch_model.py} & تولید ۲۹ لاجیت برای دسته ورودی \eng{200×200} \\
\code{test_torch_predictor.py} & اتصال پیش‌پردازش به مدل ساختگی، نگاشت خروجی و مدیریت فایل مفقود \\
\code{test_output_decoder.py} & لاجیت، احتمال، \eng{softmax}، تعداد کلاس و مقادیر نامعتبر \\
\code{test_temporal_decoder.py} & شرط دوثانیه‌ای، جلوگیری از تکرار، آزادشدن و آستانه اطمینان \\
\code{test_sequence_finalizer.py} & شرط پنج‌ثانیه‌ای، عدم نهایی‌سازی تکراری و بافر خالی \\
\code{test_label_handler.py} & رفتار حروف، فاصله، حذف، \eng{NOTHING} و برچسب ناشناخته \\
\code{test_text_buffer.py} & افزودن، حذف، فاصله، پاک‌سازی و متن نهایی \\
\code{test_mock_predictor.py} & صحت خروجی پیش‌بینی‌کننده ساختگی \\
\bottomrule
\end{tabularx}
\end{table}

تیم پروژه اجرای موفق تمام تست‌های موجود و بدون خطای بررسی \eng{Ruff} را در محیط توسعه گزارش کرده است. تنظیمات \eng{Ruff} نوت‌بوک‌ها را از تحلیل کد پایتون حذف می‌کند، زیرا خروجی و ساختار سلولی نوت‌بوک الزاماً قواعد یک ماژول تولیدی را رعایت نمی‌کند.

\section{مدیریت خطا}
موارد زیر با خطاهای اختصاصی یا پیام‌های واضح کنترل شده‌اند:
\begin{itemize}
  \item بازنشدن دوربین یا شکست در خواندن فریم
  \item نبودن فایل مدل یا ناسازگاری وزن‌ها با معماری
  \item درخواست \eng{CUDA} در سیستم فاقد \eng{CUDA}
  \item تصویر خالی، تعداد کانال اشتباه یا آمار نرمال‌سازی نامعتبر
  \item خروجی مدل با شکل نامعتبر، تعداد کلاس اشتباه یا مقادیر \eng{NaN/Inf}
  \item برچسب خالی، تکراری یا ناشناخته
  \item زمان‌های غیرصعودی در ماشین‌های حالت
\end{itemize}

\section{قابلیت تعویض اجزا}
دو کارخانه مجزا برای پیش‌بینی و گفتار وجود دارد. در مسیر پیش‌بینی، حالت \eng{mock} برای توسعه و حالت \eng{torch} برای مدل نهایی انتخاب می‌شود. در مسیر صوت نیز \eng{noop} و \eng{pyttsx3} قابل انتخاب هستند. این الگو تست‌پذیری را افزایش و وابستگی مستقیم حلقه اصلی به کتابخانه‌ها را کاهش داده است.

% ==================== Chapter 6 ====================
\chapter{تحلیل محدودیت‌ها و پیشنهادهای بهبود}

\section{شکاف دامنه میان دیتاست و وب‌کم}
داده آموزشی از محیط، زاویه و نور نسبتاً همگنی برخوردار است. عملکرد \eng{99.34\%} روی تقسیم داخلی همان منبع، دقت روی دوربین و فرد جدید را تضمین نمی‌کند. کلاس‌های بصری نزدیک مانند \eng{O/C/Q} یا \eng{V/U/W} معمولاً به زاویه انگشتان و وضوح مرزها حساس‌اند.

\section{حروف و حرکات پویا}
مدل ورودی تک‌فریم دارد؛ در نتیجه حرکت زمانی را مستقیماً مدل نمی‌کند. برخی حروف زبان اشاره مانند \eng{J} و \eng{Z} در اجرای طبیعی شامل مسیر حرکتی‌اند. مدل ممکن است یک وضعیت ثابت از این حرکت را طبقه‌بندی کند، اما برای شناخت صحیح حرکت، استفاده از دنباله فریم و معماری‌هایی مانند \eng{3D-CNN}، \eng{LSTM} یا \eng{Transformer} مناسب‌تر است.

% ==================== Chapter 7 ====================
\chapter{راهنمای اجرا و ساختار تحویل}

\section{نصب}
\begin{latin}
\begin{lstlisting}[style=pythonstyle,language=bash,numbers=none]
python -m venv .venv
source .venv/bin/activate
python -m pip install -e ".[dev,tts,model]"
\end{lstlisting}
\end{latin}

برای نصب \eng{PyTorch} سازگار با کارت گرافیک، درایور و نسخه \eng{CUDA} باید از بسته مناسب سیستم استفاده شود. در صورت نبود \eng{CUDA}، مقدار \code{device="auto"} به‌صورت خودکار \eng{CPU} را انتخاب می‌کند.

\section{اجرای کنترل کیفیت}
\begin{latin}
\begin{lstlisting}[style=pythonstyle,language=bash,numbers=none]
pytest
ruff check .
\end{lstlisting}
\end{latin}

\section{اجرای سامانه}
فایل وزن باید در مسیر زیر قرار داشته باشد:
\begin{center}
\code{models/final_model.pth}
\end{center}
سپس برنامه از ریشه پروژه اجرا می‌شود:
\begin{latin}
\begin{lstlisting}[style=pythonstyle,language=bash,numbers=none]
python main.py
\end{lstlisting}
\end{latin}

\section{ساختار فایل نهایی تحویل}
طبق دستور پروژه، تمام محتوا باید در یک فایل فشرده با الگوی نام‌گذاری تعیین‌شده قرار گیرد:
\begin{latin}
\begin{lstlisting}[style=pythonstyle,language={},numbers=none]
AI-Project3-StudentID1-StudentID2.zip
|-- main.py
|-- pyproject.toml
|-- models/
|   `-- final_model.pth
|-- notebooks/
|   |-- phase1&2.ipynb
|   `-- dataset_stats.pt
|-- reports/
|   |-- report.pdf
|   `-- report.tex
|-- src/sign_translator/
|-- tests/
|-- README.md
`-- LICENSE
\end{lstlisting}
\end{latin}

% ==================== Chapter 8 ====================
\chapter{نتیجه‌گیری}
پروژه حاضر سه بخش داده، مدل و محصول بلادرنگ را در یک زنجیره یکپارچه قرار داده است. فاز اول نشان داد داده آموزشی از نظر تعداد کلاس‌ها، اندازه تصویر و حالت رنگ ساختاری منظم دارد. در فاز دوم، شبکه چهاربخشی \eng{CNN} با افزایش داده و \eng{Dropout} به دقت آزمون \eng{99.34\%} رسید. در فاز سوم، مدل از محیط نوت‌بوک خارج و در یک برنامه واقعی مبتنی بر وب‌کم، رمزگشایی زمانی، ویرایش متن و تبدیل به گفتار به کار گرفته شد.

نقطه قوت اصلی پیاده‌سازی، جداسازی مسئولیت‌ها و توسعه افزایشی بر اساس قرارداد پیش‌بینی است. کد وب‌کم و منطق متن پیش از آماده‌شدن مدل با پیش‌بینی‌کننده ساختگی تکمیل شد و پس از تحویل وزن‌های \eng{PyTorch}، تنها لایه استنتاج جایگزین گردید. استفاده از تست‌های واحد و اعتبارسنجی صریح داده‌ها نیز ریسک خطاهای خاموش در اتصال مدل را کاهش داده است.

% ==================== Appendices ====================
\appendix
\chapter{تنظیمات نهایی سامانه}
\begin{table}[H]
\centering
\caption{مقادیر پیش‌فرض تنظیمات اجرایی}
\begin{tabularx}{\textwidth}{R{5.2cm}Y}
\toprule
\textbf{تنظیم} & \textbf{مقدار} \\
\midrule
رزولوشن دوربین & \eng{1280×720} \\
آینه‌سازی & فعال \\
اندازه ورودی مدل & \eng{200×200×3} \\
ترتیب رنگ & \eng{RGB} \\
چیدمان تنسور & \eng{NCHW} \\
مدل & \code{models/final_model.pth} \\
دستگاه & \eng{auto} (\eng{CUDA} در صورت امکان، در غیر این صورت \eng{CPU}) \\
آستانه ثبت علامت & \eng{0.80} \\
زمان تثبیت علامت & \eng{2.0 s} \\
زمان آزادسازی قفل & \eng{0.4 s} \\
زمان نهایی‌سازی & \eng{5.0 s} \\
موتور گفتار & \eng{pyttsx3} \\
کلیدهای خروج & \eng{Q} و \eng{ESC} \\
\bottomrule
\end{tabularx}
\end{table}

\chapter{قرارداد وزن‌های مدل}
فایل وزن ذخیره‌شده یک \eng{OrderedDict} شامل ۱۲ تنسور پارامتر است. ابعاد مشاهده‌شده در آزمون بارگذاری عبارت‌اند از:
\begin{latin}
\begin{lstlisting}[style=pythonstyle,language={},numbers=none]
conv1.weight                 (32, 3, 3, 3)
conv1.bias                   (32,)
conv2.weight                 (64, 32, 3, 3)
conv2.bias                   (64,)
conv3.weight                 (128, 64, 3, 3)
conv3.bias                   (128,)
conv4.weight                 (256, 128, 3, 3)
conv4.bias                   (256,)
fully_connected1.weight      (512, 36864)
fully_connected1.bias        (512,)
fully_connected2.weight      (29, 512)
fully_connected2.bias        (29,)
\end{lstlisting}
\end{latin}
این ابعاد با معماری گزارش‌شده و تعداد پارامتر کل مدل سازگارند.

\end{document}